\documentclass[10pt,a4paper]{amsart}
\usepackage[margin=19mm]{geometry}
\usepackage{amsmath,amssymb,mathtools}
\usepackage[scr=rsfs,cal=euler]{mathalfa}
\usepackage{xfrac}
\usepackage[unicode,hidelinks,bookmarks=false]{hyperref}
\newcommand{\ie}{i.e.\ }
\newcommand{\cf}{cf.\ }
\newcommand{\resp}{respectively\ }
\newcommand{\Subsection}{\subsection}
\providecommand{\nobreakdash}{\nobreak\hbox{-}}
\setlength{\emergencystretch}{2em}
\multlinegap0pt
\usepackage{amssymb}
\usepackage{cancel}
\usepackage{bm}
\usepackage{mathtools}
\usepackage{xcolor}
\usepackage{tikz}
\newtheorem{prop}{Proposition}
\usepackage{cleveref}
\crefname{prop}{Proposition}{Propositions}
\newcommand{\beq}{\begin{equation}}
\newcommand{\eeq}{\end{equation}}
\newcommand{\g}{\gamma}
\newcommand{\hx}{\hat{x}}
\newcommand{\lc}{\left(}
\newcommand{\p}{\psi}
\newcommand{\rc}{\right)}
\newcommand{\ux}{{\underline x}}
\title{Continuous-Time Heterogeneous Agent Models \\with Recursive Utility and Preference for Late Resolution}
\author{Yves Achdou, Qing Tang}
\date{Source: arXiv:2603.07782v3 (CC BY 4.0)}
\begin{document}
\maketitle

\noindent\emph{Source and licence.} Continuous-Time Heterogeneous Agent Models \\with Recursive Utility and Preference for Late Resolution. Version arXiv:2603.07782v3 (CC BY 4.0), distributed by arXiv under the Creative Commons Attribution 4.0 International licence, http://creativecommons.org/licenses/by/4.0/, as recorded in the arXiv licence metadata for this exact version and shown on the abstract page https://arxiv.org/abs/2603.07782v3. Full licence notice: \texttt{licenses/arxiv-2603.07782-CC-BY-4.0.txt}.\par\smallskip

\noindent\textit{Editorial scope.} Counted unit: the article's prop Prop:s2boundary (source0.tex lines 1050--1058). The article's complete original proof of it is delivered with it (source0.tex lines 1059--1092, one \texttt{proof} environment of 34 lines). No mathematics is written, repaired or re-derived: the statement and the proof are the article's own lines, in the article's own notation, and no retained line is substituted or re-flowed.  Retained uncounted as the source-local context the proof needs: lines 845--847; lines 881--883.  Nothing was omitted from any retained span, so the package carries no omission record.  No retained line contains a citation, so the wrapper carries no bibliography and no bibliography entry.  No retained line contains a figure environment, an image-inclusion command or drawing code, so no image is delivered and the package declares no asset dependency.  Editorial formatting only, in addition: an AMS \texttt{amsart} preamble that restates the article's own declarations and macros used by the retained lines, namely the environments \texttt{prop} on separate counters, together with the standard packages those lines need.  The retained environments print in this document as Proposition 1 in order of appearance, whereas the article prints the same items with the numbering of its own full document.  The complete native source of the article as distributed in the arXiv e-print for this exact version is bundled unchanged as \texttt{sources/195964/source0.tex}.
\begin{prop}\label{Prop:v W2}
The function $v_j$ is $W^{2,\infty}_{loc}$ in $(\ux,+\infty)$. 
\end{prop}

\begin{prop}\label{Prop: s1<0}
The optimal saving policy $s_1$ has the following properties: $s_1(x)< 0$ for all $x> \ux$ and $s_1(\ux)=0$. 
\end{prop}

\begin{prop}\label{Prop:s2boundary}
Suppose $r>0$ and 
\beq\label{rho-r}
\begin{aligned}
\lc \frac{\rho}{\theta}-r\rc(r\ux+y_2)^{-1/\psi}+\lambda_2\lc \underbrace{(r\ux+y_2)^{-1/\psi}-(r\ux+y_1)^{-1/\psi}}_{<0}\rc\geq 0.
\end{aligned}
\eeq
Then $s_2(x)<0$ for all $x>\ux$ and $s_2(\ux)=0$. 
\end{prop}

\begin{proof}
We argue by contradiction and suppose that there exists $\hx>\ux$ such that $s_2(\hx)\geq 0$. Then 
$$
c_2(\hx)\leq r\hx+y_2\quad {\text{and}}\quad Dv_2(\hx)\geq \rho(r\hx+y_2)^{-1/\psi}((1-\g)v_2(\hx))^{\frac{1-\g \psi}{\psi(1-\g)}}. 
$$
From \Cref{Prop: s1<0}, $s_1(\hx)<0$, hence
$$
c_1(\hx)>r\hx+y_1\quad {\text{and}}\quad Dv_1(\hx)<\rho(r\hx+y_1)^{-1/\psi}((1-\g)v_1(\hx))^{\frac{1-\g \psi}{\psi(1-\g)}}. 
$$

With \Cref{Prop:v W2}, if $s_2(\hx)> 0$ we differentiate the HJB equation for $v_2$ at $\hx$ to obtain
\beq\label{Eq: Dv2 2}
\lc \frac{\rho}{\theta}-r\rc Dv_2(\hx)+\lambda_2\lc Dv_2(\hx)-Dv_1(\hx)\rc=s_2(\hx)D^2v_2(\hx)+H_v(\hx,y_2,v_2(\hx),Dv_2(\hx))Dv_2(\hx)\leq 0.
\eeq
{If $s_2(\hx)=0$ we obtain \cref{Eq: Dv2 2} by differentiating in a neighborhood of $\hx$ and pass to the limit.}\par
{Let us consider the case $\frac{\rho}{\theta}-r+\lambda_2 \leq 0$. In this case,}
\beq\label{Ineq: hx}
\lc \frac{\rho}{\theta}-r\rc(r\hx+y_2)^{-1/\psi}+\lambda_2\lc \underbrace{(r\hx+y_2)^{-1/\psi}-(r\hx+y_1)^{-1/\psi}}_{<0}\rc< 0,
\eeq
hence
\beq\label{Ineq: ux}
\frac{\frac{\rho}{\theta}-r+\lambda_2}{\lambda_2}<\lc \frac{r\hx+y_2}{r\hx+y_1}\rc^{1/\p}<\lc \frac{r\ux+y_2}{r\ux+y_1}\rc^{1/\p},
\eeq
where for the last inequality we used $y_2>y_1$, $r\hx>r\ux$, $r\ux+y_j>0$ and $1/\p>0$. It is easy to see that \eqref{Ineq: ux} contradicts \eqref{rho-r}. \par
Now let us consider the case when $\frac{\rho}{\theta}-r+\lambda_2 > 0$. From \Cref{Eq: Dv2 2}, we obtain that
$$
\lc \frac{\rho}{\theta}-r+\lambda_2 \rc (r\hx+y_2)^{-1/\psi}\lc\frac{v_2(\hx)}{v_1(\hx)}\rc^{\frac{1-\g \psi}{\psi(1-\g)}}-\lambda_2 (r\hx+y_1)^{-1/\psi}<0.
$$
Since $v_1(\hx)<v_2(\hx)<0$ and $\frac{1-\g \psi}{\psi(1-\g)}<0$, we observe that
$
\frac{v_2(\hx)}{v_1(\hx)}<1$, thus
$ \lc\frac{v_2(\hx)}{v_1(\hx)}\rc^{\frac{1-\g \psi}{\psi(1-\g)}}>1$.
We recover \eqref{Ineq: hx} and then \eqref{Ineq: ux}, in contradiction with \eqref{rho-r}. We have therefore proved that $s_2(x)<0$ for all $x>\ux$ and conclude $s_2(\ux)=0$, from the state constraint condition and the continuity of $s_2$. 
\end{proof}

\end{document}
